\section{A raw residual-inversion theorem}
Let $\eta_{n,R}$ be finite measures on $\Z^3\times\R$ with characteristic functions $\Phi_{n,R}$. Frequencies are $\omega=(u,s,b)\in[-\pi,\pi]^3\times\R$, and our transform convention is $e^{i\omega\cdot z}$. Suppose there are explicit signed edge measures $E_{n,R}$ with densities $e_{n,R}$, and put $H_{n,R}=\Phi_{n,R}-\widehat E_{n,R}$. The following hypotheses are an interface, not consequences of the mechanical constructions above.

\begin{theorem}[Residual criterion for the mixed density LLT]\label{thm:LLT}
Assume the following uniformly for $R\in I$.

\smallskip\noindent
(i) There are $c,C,\delta>0$, $\rho<1$, means $m_R\in\R^4$, and symmetric matrices $cI\le\Sigma_R\le CI$ such that, for $|\omega|\le\delta$,
\[
 \Phi_{n,R}(\omega)=e^{inm_R\cdot\omega}
 \{A_R(\omega)\lambda_R(\omega)^n+O(\rho^n)\},
\]
where $A_R(0)=1$, $|A_R(\omega)-1|\le C|\omega|$,
$\log\lambda_R(\omega)=-\tfrac12\omega^T\Sigma_R\omega+O(|\omega|^3)$, and $|\lambda_R(\omega)|\le e^{-c|\omega|^2}$.

\smallskip\noindent
(ii) $H_{n,R}\in L^1(\T^3\times\R)$ and
\begin{equation}\label{eq:tailcriterion}
 \int_{|\omega|>n^{-2/5}}|H_{n,R}(\omega)|\dd\omega=o(n^{-2}).
\end{equation}

\smallskip\noindent
(iii) $\|E_{n,R}\|_{\TV}=o(n^{-2})$ and, for each fixed $K$,
\[
 n^2\mathop{\rm ess\,sup}_{|(z-nm_R)/\sqrt n|\le K}|e_{n,R}(z)|\longrightarrow0.
\]
Then $\eta_{n,R}$ has a density with respect to counting measure times Lebesgue measure, and on each such central window
\begin{equation}\label{eq:LLT}
 n^2 p_{n,R}(k,m,t)=
 \frac{\exp(-\xi^T\Sigma_R^{-1}\xi/2)}{(2\pi)^2\sqrt{\det\Sigma_R}}+o(1),
 \quad \xi=\frac{(k_1,k_2,m,t)-nm_R}{\sqrt n},
\end{equation}
with uniform essential-supremum error. No smoothing of $\eta_{n,R}$ is used.
\end{theorem}
\begin{proof}
Fourier inversion of $H_{n,R}$ gives a continuous density of the signed measure $\eta_{n,R}-E_{n,R}$. To justify this without assuming a density, take each lattice Fourier coefficient of $H$. It is an integrable roof transform of the corresponding signed component. Convolve that component with a Gaussian approximate identity, use integrability to pass to the inverse-transform limit, and test against compactly supported continuous functions. Uniqueness of finite Fourier transforms identifies each component, hence the full measure. This is a proof of inversion, not a change to the datum. Adding the known density of $E_{n,R}$ gives $p_{n,R}$.

On $|\omega|\le n^{-2/5}$, replacing $H$ by $\Phi$ costs at most a fixed volume times $\|E\|_{\TV}=o(n^{-2})$. Inserting (i), let $v=\sqrt n\,\omega$. The cubic error is uniformly $O(n^{-1/5})$ on this region; the bound on $\lambda$ supplies an integrable Gaussian dominator. Uniform Taylor bounds, truncation first to $|v|\le M$, and then $M\to\infty$, give the inverse Gaussian with error $o(n^{-2})$. The inverse normalization is $(2\pi)^{-4}$ and the Gaussian integral is $(2\pi)^2/\sqrt{\det\Sigma_R}$, giving the constant in \eqref{eq:LLT}. Hypothesis (ii) bounds the remaining frequencies, and (iii) bounds the added density in the central window. All estimates are uniform in $\xi$ in that window.
\end{proof}

For weighted inserted measures, the same proof applies with $A_R(0)=a_R$ and uniformly bounded amplitudes, provided the edge and tail hypotheses are checked for those measures as well. It gives $a_R$ times the Gaussian. It does not assert uniformity for arbitrary $n$-dependent path functionals without such estimates.

\begin{lemma}[An explicit frequency splice]\label{lem:splice}
Suppose a residual transform has a compact minor-arc bound $Ce^{-cn}$ and, for $1\le|b|\le e^{\kappa n}$, a bound $C(1+|b|)^A e^{-cn}$. Beyond $e^{\kappa n}$ suppose it is bounded by
\[
 C_M\rho^n(1+|b|)^{-2}+C_Me^{c_M n}(1+|b|)^{-M}.
\]
If one fixed integer $M>1$ and one $\kappa>0$ satisfy
\begin{equation}\label{eq:splice}
       \frac{c_M}{M-1}<\kappa<\frac{c}{A+1},
\end{equation}
these regions contribute exponentially small $L^1$ mass. Together with the Gaussian major-arc bound between $n^{-2/5}$ and $\delta$, this proves \eqref{eq:tailcriterion}.
\end{lemma}
\begin{proof}
The middle integral is bounded by $C\exp(-(c-(A+1)\kappa)n)$. The two outer integrals are bounded by $C_M\exp((\log\rho-\kappa)n)$ and $C_M\exp(-( (M-1)\kappa-c_M)n)/(M-1)$. The torus has finite volume. The major-arc Gaussian tail is bounded by a polynomial factor times $e^{-c n^{1/5}}$, and the compact remainder is exponential. Strict positivity of the three exponent margins proves the claim.
\end{proof}
The inequality \eqref{eq:splice} must be verified with the actual derivative-growth constants. The phrases ``take $\kappa$ small'' and ``then take $M$ large'' do not imply it: for example, $c_M=2(M-1)$ and $c/(A+1)=1$ leave no possible $\kappa$. This checks a proposed estimate, not the truth or falsity of a central LLT. The residual formulation also requires an independently proved bound for the sum of all extracted edge terms; Proposition~\ref{prop:subtract} treats one physical family only.

